\chapter{Селектор совместимой комплексной структуры и положительной метрики}
\label{ch:v8-horizontal-phase-cotangent-complex-structure-metric-selector-gate}

Котангенциальный родитель дал ненулевое отображение момента, но сохранил
волоконную фазу. Следующий возможный источник различения --- совместимая
комплексная структура $J$ и положительная форма
\begin{equation}
 J^2=-I,qquad g_J(u,v)=\Omega(u,Jv).
 \label{eq:v8-cotangent-compatible-complex-structure}
\end{equation}
Само существование таких $J$ стандартно, но их выбор требует метрических
данных. Проверим, достаточно ли уже построенной следовой метрики проекта.

Здесь возникает типовой барьер, который не был виден в абстрактном
симплектическом гейте. Вещественная размерность нового носителя равна $52$,
тогда как полный следовой полево-шумовой репер имеет размерность $42$:
\begin{equation}
 \dim_{\mathbb R}\mathcal T_{26}^{\mathbb R}=52,qquad
 \dim_{\mathbb R}\mathcal F_{\rm tr}=42,qquad 52-42=10.
 \label{eq:v8-cotangent-trace-carrier-mismatch}
\end{equation}
После отбеливания невырожденной старой следовой формы можно считать её
матрицей $I_{42}$. Для любого отображения
$P:\mathbb R^{52}\to\mathbb R^{42}$ её обратный образ удовлетворяет
\begin{equation}
 G_P=P^{\mathsf T}P,qquad
 \rank G_P\leq42,qquad \dim\ker G_P\geq10.
 \label{eq:v8-cotangent-trace-pullback-bound}
\end{equation}
Точный представитель $P=(I_{42}\;0)$ достигает границы: его ранг равен
$42$, а нульность --- $10$. Если использовать лишь старую тридцатимерную
вещественную метрику переноса, граница ещё сильнее:
\begin{equation}
 \rank G_{\rm transfer}=30,qquad
 \dim\ker G_{\rm transfer}=22.
 \label{eq:v8-cotangent-transfer-pullback-bound}
\end{equation}
Следовательно, существующая следовая форма не является положительной
метрикой на всём новом носителе.

Более того, даже в наиболее благоприятном случае, когда старые $42$
направления образуют симплектическое подпространство, продолжение не
единственно. Возьмём
\begin{equation}
 \Omega_{52}=J_2^{\oplus26},qquad
 J_2=\begin{pmatrix}0&1\\-1&0\end{pmatrix},qquad
 G_s=I_{42}\oplus
 \bigoplus_{a=1}^{5}\begin{pmatrix}s&0\\0&s^{-1}\end{pmatrix},quad s>0.
 \label{eq:v8-cotangent-compatible-metric-family}
\end{equation}
Положим
\begin{equation}
 \mathcal J_s=-\Omega_{52}G_s.
 \label{eq:v8-cotangent-compatible-J-family}
\end{equation}
Тогда точная блочная проверка даёт
\begin{equation}
 \mathcal J_s^2=-I_{52},qquad
 \Omega_{52}\mathcal J_s=G_s>0,qquad
 P G_sP^{\mathsf T}=I_{42},qquad \det G_s=1.
 \label{eq:v8-cotangent-compatible-family-identities}
\end{equation}
В частности, $s=1$ и $s=2$ дают две различные положительные совместимые
метрики и две различные комплексные структуры, которые имеют одно и то же
ограничение на весь старый следовой носитель.

Это точно воспроизводит раннюю проблему Тома III, где одна каноническая
пара допускала семейство $J_s$. Новый уровень увеличил носитель и дал
конкретное отображение момента, но не породил десять недостающих
метрических направлений. Общий факт о непустом пространстве совместимых
комплексных структур не предоставляет селектора конкретного элемента
\cite{v8-compatible-metric-da-silva,v8-compatible-metric-guillemin}.

\begin{theorem}[Дефицит следового метрического продолжения]
Никакой обратный образ существующей $42$-мерной следовой метрики не может
быть положительно определён на $52$-мерном симплектическом носителе: его
ядро имеет размерность не меньше десяти. Даже после произвольного
положительного заполнения ядра остаётся непрерывное семейство совместимых
метрик и комплексных структур с одинаковым старым ограничением.
\end{theorem}

\begin{nogocriterion}[No-go следового селектора]
Нельзя объявлять $J=-\Omega$ каноническим только после выбора единичной
метрики на новых направлениях: эта единичность является дополнительным
правилом. Существующий след фиксирует не более $42$ из $52$ вещественных
координат и не различает семейство $G_s$. Поэтому комплексная структура,
положительный родитель и масса горизонтальной фазы ещё не выведены.
\end{nogocriterion}

\begin{protocol}
\begin{enumerate}
 \item новый симплектический носитель: \textbf{$52$ вещественных};
 \item полный старый следовой репер: \textbf{$42$ вещественных};
 \item минимальный дефицит метрики: \textbf{$10$};
 \item дефицит transfer-only метрики: \textbf{$22$};
 \item положительный обратный образ старой метрики: \textbf{невозможен};
 \item совместимые положительные продолжения: \textbf{существуют};
 \item два различных точных продолжения: \textbf{$s=1,2$};
 \item старое следовое ограничение у них совпадает: \textbf{да};
 \item единственная комплексная структура: \textbf{не выбрана};
 \item подъём горизонтальной фазы: \textbf{не получен};
 \item следующий гейт: \textbf{происхождение метрики на десяти недостающих
 направлениях}.
\end{enumerate}
\end{protocol}

Машинный сертификат сохранён в
\path{s2t_v8_horizontal_phase_cotangent_complex_structure_metric_selector_gate_results.json}.

\begin{thebibliography}{9}
\bibitem{v8-compatible-metric-da-silva}
A.~Cannas da Silva,
\emph{Lectures on Symplectic Geometry}, Springer, 2001;
arXiv:math/0505366.
\bibitem{v8-compatible-metric-guillemin}
M.~Abreu,
\emph{K\"ahler geometry of toric manifolds in symplectic coordinates},
in Symplectic and Contact Topology, Fields Inst. Commun. 35 (2003);
arXiv:0912.0491.
\end{thebibliography}